

%%%% Chapter formatting %%%%
\titleformat{\chapter}[display]
    {\normalfont\bfseries}
    {\fontsize{24}{28}\selectfont Chapter \thechapter} {1.5cm} % Space between
    {\fontsize{36}{42}\selectfont}

\titlespacing*{\chapter}
    {0pt}   % Left margin
    {1cm}   % Space before chapter
    {1.5cm}   % Space between title and text

%%%% Page design %%%%
\setlength{\parindent}{0pt} % remove paragraph indentation
\setlength{\parskip}{1em}   % add vertical space between paragraphs


%%%% Custom Comment %%%%%
% uses pdfcomment
\newcommand{\showcomment}[3][]{%
  \pdfmarkupcomment[%
    markup=Underline,       % underline the text
    color=blue,            % highlight color
    author={#1},             % optional author name
    opacity=0.4,             % how strong the yellow looks (0–1)
    subject={Comment},       % shows up as title in popup
    open=false               % popup starts closed (click to open)
  ]{#3}{#2}%
}
% example:
% \showcomment[Lukas]{The comment}{The text you want to comment}

%%%% Auto Chapter %%%%
\newcommand{\inputchapter}[2]{%
    \def\basefolder{chapters}%
    \immediate\write18{ls -1v \basefolder/#1/#2.*.tex > filelist.tmp}%
    \newread\filelist
    \openin\filelist=filelist.tmp
    \loop
        \read\filelist to \filename
        \unless\ifeof\filelist
        \input{\filename}%
    \repeat
    \closein\filelist
}


%%%% Auto Preface %%%%
\ExplSyntaxOn
\seq_new:N \l__group_members_seq
\int_new:N \l__group_member_count_int

% Define one member
\NewDocumentCommand{\member}{mm}
{
    \seq_put_right:Nn \l__group_members_seq
    {
        \begin{minipage}[b]{0.45\textwidth}
            \centering
            \rule{\textwidth}{0.45pt}\\
            #1\\
            {\footnotesize #2}
        \end{minipage}
    }
}

% Define the complete member section
\NewDocumentEnvironment{members}{}{\seq_clear:N \l__group_members_seq}
{
    \int_set:Nn
        \l__group_member_count_int
        {\seq_count:N \l__group_members_seq}
    \vspace{3em}
    \vspace{3cm}
    \vfill
    \noindent

    % Complete pairs
    \int_step_inline:nn
    {
        \int_div_truncate:nn
            {\l__group_member_count_int}
            {2}
    }
    {
        \seq_item:Nn
            \l__group_members_seq
            {\int_eval:n {2 * ##1 - 1}}
        \hfill
        \seq_item:Nn
            \l__group_members_seq
            {\int_eval:n {2 * ##1}}

        \par
        \vspace{3\baselineskip}
    }

    % Remaining member if the number is odd
    \int_if_odd:nT {\l__group_member_count_int}
    {
        \begin{center}
            \seq_item:Nn
                \l__group_members_seq
                {\l__group_member_count_int}
        \end{center}
    }
}

\ExplSyntaxOff